\documentclass[a4paper,12pt]{article}

% Package Imports
\usepackage{styles/ibu-thesis}
\usepackage{subcaption}
\usepackage{graphicx}
\usepackage{float}
\usepackage[backend=biber,style=numeric]{biblatex}
\addbibresource{sample.bib}
% Macros
\newcommand{\writecommand}[2]{
    \begin{center}
        \texttt{\textbackslash #1\{\##2\#\}} \\
    \end{center}
}

% Thesis Title
\thesistitle{CALORIE CRAFTER}

% Author Information
\authors[1]{OZAN ŞENTÜRK}
\authorsid[1]{119200080}
\authors[2]{ALPER HAKAN BAŞER}
\authorsid[2]{116200052}
\authors[3]{ERCÜMEND KAYHAN}
\authorsid[3]{120200060}

% Supervisor Information
\supervisors[1]{Dr. MURAT ORHUN}

% Affiliation Information
\faculty{Faculty of Engineering and Natural Sciences}
\degree{Bachelor of Science}
\department{Department of Computer Engineering}
\submissiondate{June, 2024}

\abstract{
Calorie Crafter website offers the proper guidance to consumers in terms of how many calories and which nutrients they should take daily. Incorporating web 2.0 technologies and genetic selection, it has food guidelines compatible with specific nutrient requirements and life style profiles. Concerning friendly assimilation, the design concentrates on functionality to encourage the consumption of the correct portions of foods. This document aims to present the UI design and development of this report by exploring the general skills in UI design, UX optimization, and risk analysis. IEEE, IET, EU, Turkish standards, cultural identity, reliability, and safety provide high quality in this project. Those are such fragments as daily menu options, the possibility of creating a meal, as well as the section of the user’s profile. The social benefit of the project focuses on the general well-being of individuals while the environmental impacts focus on encouraging healthier diets hence impacting positively on health, prompting people to adopt sustainable diets hence drawing a positive impact on health and the economy by providing a cheaper solution to dietary management. With the help of the developed Calorie Crafter web-site, it is possible to appreciate the efforts of the team for creating a new outlook on the personal approach to the people in the framework of the nutrition process in order to have a healthier life.
}

\begin{document}

% Preliminaries
\maketitle

\addcontentsline{toc}{section}{\abstractname}
\makeabstract

\addcontentsline{toc}{section}{\contentsname}
\maketableofcontents

\newpage
\pagenumbering{arabic} % Do not change
\setcounter{page}{1} % Do not change

% Sections
\section{Introduction}
{\small

As for the present world with a focus on health and exercising, the controlling of diet of a person is very crucial. This is why the Calorie Crafter website is designed to offer personal meal plans of incorporation the recommended daily calories depending on users’ preferences, activity levels, and essential nutritional needs. Using modern web technologies and genetic algorithms to deliver individuals with the optimal diets to meet their targeted fitness goals in the most convenient and effective way possible.

The primary goal of the Calorie Crafter is a concept of the website where it will easily help the user track his/her food and beverage intake in terms of calorie and nutrient contents. Thus, such a navigation of detailed meal recommendations for each day, the possibility to custom and create one’s own meals, as well as complete account profile management will help inform people and provide them with the tools necessary to make correct food choices This tool is intentionally conceived to be user-friendly so basic functions and features including recommendation, sustaining meal plan creation, and customization are easy to find and operate.

As this vision suggests, this project was bound to be implemented by a group of hard-working students who worked collectively toward the realization of this vision. Both team members did their best by adhering to the consolidated skills and knowledge that were applied to ensure the high quality of the project and timely completion of the job. Parts of the development process included strict engineering practices, such as IEEE, IET, EU, and Turkey standard practices in order to provide a quality, safe, and stable website for a cervical cancer awareness organization.

Therefore in this report, there is a focus on analyzing the process of designing and developing the Calorie Crafter website where aspects to be covered include aspects like the UI design, UX development risks in developing the site, and the engineering requirements that need to be met. We will also briefly expound on the social, environmental, as well as economic effects of the project, in addition to outlining the inputs from the members of the project team.

The current state of the Calorie Crafter website should be viewed as a testament to the team’s commitment to creating a meaningful utility for supporting nutrition. Thus, emphasizing a user-oriented approach, compliance with high standards, and highlighting the issues now on the project can be considered as the one that has really created a platform that can help people achieve their goals in using the proper nutrition depending on their desires. The outlined process of the comprehensive development and the identified advantages of the project serve to underscore the role of individualized nutrition in the contemporary context of the health sector.


\section{Design}
The layout of the Calorie Crafter core website to the actual presentation to the users follows a functioning and visually exquisite creation that can effectively guide the users through the different processes of tracking daily intake of calories and nutrients. Their layout is simple to enable a user to move within different parts of the website in one click depending on what they are in need off; for instance, meal recommendations, creating new meals, or modifying profiles.

\subsection{User Interface (UI)}
\subsubsection{Dashboard and Daily Meal Recommendation}
In the main area of the app, we provide users with meal suggestions to suit their body values and goals. By clicking on the Generate button, users will be able to view our unique daily menu suggestions and if they don't like the menu, they can change it by pressing the button again. The quantities and values of our suggestions can be seen in the pictures of the dishes. In the section below our recommendations, users can view the total Calorie, protein, carbohydrate, and fat values of the daily menu we recommend for them. On the right side of the page, 5 random dishes from our system are displayed. These meals change every time the page is refreshed. On the left side of our page, you can see our dashboard. On our dashboard, there are buttons to access our home page, create meal page and profile page.
}
\begin{figure}[h]
    \centering
     \fbox{\includegraphics[width=10cm]{mainpage.jpeg}}
    \caption{Daily Menu Recommendation and Dashboard}
    \label{fig:example}
\end{figure}
\newpage
\subsubsection{Create a Meal Tab}
On the “Create e Meal” page, users will be able to save the meals they have eaten throughout the day, view the calorie, carbohydrate, protein, and fat values they have taken, and track them in comparison with the daily values calculated by our algorithm for them thanks to the progress bars. There is also a date selection screen on the right side of the page above the progress bars. Users will be able to view the values they have taken on the selected date there, as well as the meals they have eaten during the day on the selected date on the right side of the page. For the convenience of users, we have filtering options under the name Breakfast-Lunch-Dinner and search options with the search bar.


\begin{figure}[h]
    \centering
     \fbox{\includegraphics[width=10cm]{create2.jpeg}}
    \caption{Create a Meal}
    \label{fig:example}
\end{figure}
\subsubsection{Profile Management Tab}
On our profile page, users can view current height, weight and BMI values.If users want to change height, weight, activity status or Goal Factors, they can update these values on our profile page. According to these updated values, the BMI value is calculated and updated again . According to the updated values, the daily calorie, protein, carbohydrate, and fat values that the user should take are calculated and changed again by our algorithm.
\begin{figure}[h]
    \centering
    \includegraphics[width=10cm]{profile.jpeg}
    \caption{Profile}
    \label{fig:example}
\end{figure}
\newpage
\subsection{User Experience (UX)}
\subsubsection{Visual Design}
The layout of the design of the website of Calorie Crafter is simple and has a contemporary look. The chosen colors are health-related and calm colors. To attract the visitors’ attention and make the interface more attractive, high-quality images of meals and fitness activities are applied.

\subsubsection{Accessibility}
Navigability is one of the ways through which the Calorie Crafter website has been approached. Corrective measures have been taken, as regards the website’s’ accessibility to persons with disabilities. This includes:
Implementing the use of alt text for the images as a way of helping the screen readers.
Meeting the requirements for the contrast of colors between the text and the background.
Keyboard access to all the content and interface controls of the website.

The layout of the Calorie Crafter website is very attractive and developed exclusively for customer interaction. Due to simple and clean design, clear and concise presentation of nutritional information, and easy-to-use functions and features, the website successfully assists users in achieving their dietary objectives and leading a healthy life.




\section{Risk Analysis}
\subsection{Potential Risks}
\subsubsection{Algorithm Performance}
An emergent risk was in relation to the effectiveness of the genetic algorithm when designing the specific diet plans on the basis of the given keys. We had to experiment with several little values that control each GA, including the mutation rate, crossover rate, as well as selection criteria to offer the best solution. There was, however, a challenge in how to prevent the algorithm from converging too early or getting stuck at the local extremum, which in the case of dieting would have resulted in presenting our users with a suboptimal diet plan.

\subsubsection{Data Quality and Availability}
 Problems were encountered with respect to the quality and accessibility of users’ data which in one way  another affected our genetic algorithm considerably. This meant that the data inputs for food choices, medical history of individuals, including exercise regimens, should be correct and also encompass all necessary details that audience provides. Finally, collecting the genetic data was again a privacy issue, thus we had to follow the rules set by GDPR or HIPAA among others.

\subsubsection{User Adherence and Satisfaction}
The second important risk we can distinguish can be linked to the non-compliance of users with the recommended diet plans or the satisfaction with the website. However, the algorithm provides individual recommendations for the diet, and some users can have challenges following the prescribed diets to the letter daily. Hence, we pay attention to these factors in designing the website and making users motivated more adherent, and satisfied.

\subsubsection{Health Risks and Liability}
Recommending individual diets based on genetic predispositions introduced inherent health risks, and we had to be cautious about liabilities. It was possible that some of the recommended diets could be capable of causing detrimental health consequences where there would be worsening of pre-existing health complications. This risk was why it was essential to ensure that the diet plans recommended were safe and would be effective for the patients.

\subsubsection{Algorithm Bias and Fairness}
We also struggled with how the proposed algorithms can be biased in what they recommend for diets. The potential issue with our genetic algorithm is that it might incorporate bias related to ethnicity or other aspects entrenched in the datasets used. Some biases that are relevant to dieting were however taken in advance to give equal chances to all in order to remove bias.

\subsubsection{Competitive Landscape and Market Acceptance}
Last but not the least, there emerged the issue of competition and the need to address market acceptance for the established diet website. The various other platforms provided other strategies on how to reach the principle of personalized nutrition, which we had to disentangle skillfully. It was important to make sure the users trusted our way of doing things and meeting their relevant needs and issues in the course of the project.

\subsection{Risk Management}
\subsubsection{Algorithm Performance}

\textbf{Contingency Plan:} We have developed a robust testing framework that allowed us to evaluate the performance of the genetic algorithm under various scenarios. In instances where the algorithm converged prematurely or produced suboptimal results, we successfully implemented fallback mechanisms to revert to previous iterations or apply corrective measures.

\noindent\textbf{Regular Monitoring:} We established a system for continuous monitoring of the algorithm's performance. By tracking key performance indicators such as convergence rate, diversity of solutions, and fitness scores, we were able to identify deviations from expected behavior and address them proactively.

\noindent\textbf{Alternative Approaches:} We explored alternative optimization techniques such as simulated annealing or particle swarm optimization to complement our genetic algorithm. This ensured that we had backup strategies in place to achieve our objective of generating personalized diet plans effectively.
\subsubsection{Data Quality and Availability}

\textbf{Data Validation:} We implemented rigorous validation procedures to ensure the quality and consistency of user-provided data. Through automated data validation scripts and manual verification processes, we successfully mitigated the risk of poor data quality.

\noindent\textbf{Privacy Compliance:} We strictly adhered to data protection regulations such as GDPR or HIPAA to safeguard user privacy. By implementing encryption protocols, access controls, and anonymization techniques, we ensured compliance with privacy regulations.

\noindent\textbf{Data Augmentation:} In cases where user data was insufficient, we employed data augmentation techniques such as synthetic data generation and leveraging publicly available datasets. This allowed us to supplement our training data effectively.
\subsubsection{User Adherence and Satisfaction}

\textbf{Behavioral Nudges:} We integrated behavioral psychology principles into the design of our website to encourage user adherence to recommended diet plans. Through personalized reminders, progress-tracking features, and social support mechanisms, we successfully fostered accountability and motivation among users

\noindent\textbf{Feedback Mechanisms:} We established regular feedback loops to gather user input and assess satisfaction levels. Through surveys, user reviews, and analytics tracking, we identified areas for improvement and addressed user concerns promptly.
\subsubsection{Health Risks and Liability}

\textbf{Expert Consultation:} Thus, to help us avoid some health risks that are associated with diets that we recommend to our clients, we worked with nutritionists, dieticians, and health care providers. They confirmed compliance of the diet plans with the internationally accepted guidelines thus embracing safety and efficiency in our diet plans.
\subsubsection{Algorithm Bias and Fairness}

\textbf{Bias Detection and Mitigation:} To ensure that our results are not being skewed by an unintentional bias, we introduced programs that would help to prevent bias in our genetic algorithm. I thus performed fairness audits and sensitivity analyses to check whether demographic data influenced its dietary plans and refined the algorithm of the AI accordingly.

\noindent\textbf{Diverse Training Data:} Bring into focus that training data was fairly collected over the target population in order to avoid bias towards specific groups. There was no evidence of algorithmic biases in the Midwestern population by actively seeking samples from these groups and by adding fairness constraints to the objective function.
\subsubsection{Competitive Landscape and Market Acceptance}

\textbf{Market Research:} To ensure that we were accurate in our product range and package offerings to the market we undertook daily market research to identify the existing competitors and changing customer needs. This allowed us to be responsive in managing our strategies and to be able to position ourselves well in an environment that can be rather saturated in many cases.

\noindent\textbf{Value Proposition:} Informative communication, focusing on the original value proposition of the diet website, our key advantage – the genetic algorithm, and the“call for action” based on its practical application allowed for improving the market acceptance and brand credibility.

\section{Standards}
In designing and implementing our personalized nutrition and diet recommendation website, it is essential to adhere to relevant engineering standards to ensure quality, safety, and reliability. This includes compliance with standards from IEEE, IET, EU, and Turkish regulations, as well as adherence to the engineering code of conduct. Below are the key standards and guidelines considered for this project
\subsection{IEEE Standards}
\subsubsection{Software Engineering - Software Life Cycle Processes Maintenance (IEEE 14764-2006)}

This standard outlines the processes involved in maintaining software throughout its lifecycle. Our project adheres to this standard by implementing systematic procedures for software updates, bug fixes, and enhancements to ensure continuous improvement and reliability. \cite{ieee14764}.


\subsubsection{Recommended Practice for Software Requirements Specifications (IEEE 830-1998)}
Ensuring clear and detailed software requirements is crucial for the success of our project. This standard provides guidelines for creating comprehensive software requirements specifications (SRS) that describe the functionality, performance, and constraints of the software system. \cite{ieee830}


\subsubsection{System and Software Verification and Validation (IEEE 1012-2016)}
System and Software Verification and Validation: This standard provides guidelines for the verification and validation of software systems to ensure that they meet specified requirements and are fit for their intended use. Our project follows this standard to ensure the accuracy and effectiveness of the genetic algorithm used for personalized diet recommendations.\cite{ieee1012}


\subsection{IET Standards}
\subsubsection{IET Code of Practice for Cyber Security and Safety}
As our project handles sensitive user data, including dietary preferences, medical history, and lifestyle factors, adhering to cyber security and safety standards is paramount. This code of practice helps ensure that data is securely stored and transmitted, protecting user privacy and complying with data protection regulations.\cite{iet2020}


\subsubsection{IET Guide to the Engineering of Secure Systems}
This guide provides best practices for the design and implementation of secure systems. By following these guidelines, we ensure that our system is robust against cyber threats and vulnerabilities, maintaining user trust and data integrity. \cite{iet2016}


\subsection{EU Standards}
\subsubsection{GDPR (General Data Protection Regulation) (Regulation (EU) 2016/679 - GDPR)}
This regulation governs data protection and privacy for individuals within the European Union. Our project complies with GDPR by implementing strict data handling and storage practices, ensuring user consent, and providing users with control over their personal data. \cite{gdpr2016}


\subsubsection{Medical Device Software—Software Life Cycle Processes (EN 62304:2006)}
Given the potential health impacts of our personalized diet recommendations, this standard, which outlines the life cycle requirements for medical device software, is relevant. Although our project is not classified as a medical device, adhering to these principles helps ensure that our software is developed and maintained to a high standard of quality and safety.\cite{iso62304}


\subsection{Turkish Standards}
\subsubsection{Information Security Management Systems (ISMS) (TS ISO/IEC 27001)}
This standard specifies the requirements for establishing, implementing, maintaining, and continually improving an information security management system. Compliance with TS ISO/IEC 27001 ensures that our project follows best practices in information security, protecting user data from breaches and unauthorized access. \cite{tse27001}


\subsubsection{Quality Management Systems (TS EN ISO 9001)}
This standard outlines the criteria for a quality management system and is used to ensure that organizations can consistently provide products and services that meet customer and regulatory requirements. By adhering to TS EN ISO 9001, our project emphasizes quality assurance and continuous improvement. \cite{tse9001}


\section{Conclusion}
\subsection{Realistic Constraints}
Several factors have to be taken into consideration in designing the Calorie Crafter website and there are some realistic constraints that have been considered as follows in designing the ideal Calorie Crafter website. Among them, it involves defining a problem and then finding a solution that is an engineering project, and the solution has to follow specifications and standards, the resources available and the budget or time limit available. In addition, the capabilities of the engineering solutions on social, economic, and ecological factors were also evaluated.
\subsubsection{Budget Constraints}
The proposed project was sketched in order to run in a small scale since the Belkat used limited funds and other resources to initiate the projects. Better solutions for the development, design, and testing of the software were very much needed which are the thing a practitioners wants. Low-cost strategies were implemented with an eye on cost outlay with quality maintained at par with and very close to that of standard open source technologies and frameworks were applied.

\subsubsection{Timeline Constraints}
There was also a respect for time in the completion of the project with a set timeframe laid out to complete the project. These methods were adopted to ensure that activities and schedules were well coordinated as project development cycles were implemented. Schedules were established at intervals and the projects were held accountable to them to ensure that they would not stray far off schedule.

\subsubsection{Engineering Standards}
Knowledge of good engineering practices defined in IEEE and IET as well as EU regulations and Turkish standards was necessary in order to ensure high quality, safety and reliability of the website. Procedures including ISO 14764-2006 on Software Maintenance and General Data Protection Regulation (GDPR) were adhered to.



\subsection{Social, Environmental, and Economic Impact}
From the uncollected information, it is seen that the website of Calorie Crafter is designed with the ultimate goal of assisting consumers in improving their nutrition and fitness. When analyzing the merits and demerits of this project, social, environmental, and economic impacts can in no uncertain terms be described as complex and far-reaching in various ways.
\subsubsection{Social Impact}
The stand-out social effect of the Calorie Crafter website is its contribution to enhancing the health of society. It offers better and closer to real life and the particular needs of an individual, nutrition consulting, enabling him or her to make the right choices that will enhance a better and healthier life. \textbf{This has several social benefits:} 
\vspace{1cm}
 \\ \noindent\textbf{Improved Health and Well-being:} Users also get the ability to manage their diets which is crucial as it helps lessen the number of cases of diseases related to poor diet thus incidences of obesity diabetes and cardiovascular diseases. This can be attributed to the fact that the communities’ health status is generally improved.

\noindent\textbf{Education and Awareness:} Users are aware of the nutrition profile of the foods and general diet that is healthy for the human body. It may result in health conscious eating behaviors, which can create a positive change, there by increasing awareness.

\noindent\textbf{Accessibility:} By enabling clients to receive recommendations of affordable or free of charge, Calorie Crafter ensures that it can serve clients that can be unable to seek advice from dietitians or nutritionists.


\subsubsection{Environmental Impact}
While the primary focus of the Calorie Crafter website is on nutrition, there are indirect environmental benefits associated with the project:\\
\noindent\textbf{Promotion of Sustainable Eating:} Therefore, when the app is capable of delivering all the detailed nutritional information, the chances are high that the users of the platform will make better decisions of choosing organic foods and taking plant-based meals among other positive changes. This can in turn result in a reduced environmental cost in the production as well as the processing of food products.

\noindent\textbf{Reduction in Paper Waste:} In line with this the use of the digital platform means that there shall be no need to use printed copies of the diet plans and guides hence; the main benefits include; lesser usage of paper or less environmental pollution.

\subsubsection{Economic Impact}
The economic impact of the Calorie Crafter website is notable in several areas:

\noindent\textbf{Cost Savings for Users:} Individuals can easily get their nutritional plans from their online sources and without needing to spend large amounts of money eons recommended dietitian and nutritionists. Coupled with the fact that it can be used consistently for dietary health, it also makes this a cost-effective intervention.

\noindent\textbf{Healthcare Cost Reduction:} As society becomes more conscious of their dietary habits; this means that incidences of diseases related to diets would be less frequent, hence the costs both on individuals and the society would go down. Counseling actually entails preventative measures toward major health complications and this in the long run has major implications for cost reduction.

\noindent\textbf{Market Opportunities:} The implementation of the project can create commercial potential for relations with the outlet’s food brands, fitness center, and healthcare providers implying economic opportunities attached to the related industry. Also, there have to be innovative ways to monetize the platform, such as incorporating a fee for special updates or specific coaching packages.

In conclusion, the Calorie Crafter website, while primarily focused on providing personalized dietary recommendations, has significant social, environmental, and economic impacts. It promotes better health and well-being, encourages sustainable eating habits, and offers a cost-effective solution for dietary management.


\subsection{Cost Analysis}
Conducting a comprehensive cost analysis for the Calorie Crafter website project involves evaluating the expenses associated with labor, design, and the various technical components used in the development process. Below is a detailed breakdown of these costs:

\subsubsection{Labor Costs}
The project was a group affair and it was done by three members of the students. In total, students spent 120 hours on the given project, so it’s possible to calculate hourly rate per student at \$40.
\\ \noindent\textbf{Total Hours Worked:} 120 contact hours

\noindent\textbf{Hourly Rate:} \$10

\noindent\textbf{Total Labor Cost:} 120 hours * 10 = the amount that is 1200.


\subsubsection{Design Costs}
Design costs are constituted by acquiring a diet website design template.

\noindent\textbf{Diet Website Design Template:} Rather it can range from \$25 on the lower end up to \$50 for more complex treatments.
 \\ \noindent\textbf{Total Design Cost:} \$50

\subsubsection{Technical Costs}

Technical costs are the costs of a Website Design Template for the website and the registering of its domain name.

\noindent\textbf{Website Design Template:} \$50

\noindent\textbf{Domain Registration:} \$20

\noindent\textbf{Total Technical Cost:} \$70

\newpage
\subsubsection{Summary of Costs}
\noindent\textbf{Labor Costs:} \$1,200

\noindent\textbf{Design Costs:} \$50

\noindent\textbf{Technical Costs:} \$70

\noindent\textbf{Total Project Cost:} \$1,420
\vspace{1cm}

\noindent\Summing Up the costs of website development, which include labor, design, and technical details, the Calorie Crafter cost \$ four, nine hundred & twenty of money. This proposed detailed cost breakdown outlines the investment necessary to offer a platform of great quality and implementation of a user-first approach. These costs need to be well controlled and all utilization of the resources should be done parsimoniously in order to offer the target users a worthy value addition solution within the available project budget.


\subsection{Contributions}
The successful development of the Calorie Crafter website is the result of collaborative efforts and contributions from all team members. Each member played a crucial role in various aspects of the project, ensuring that all tasks were completed efficiently and to a high standard. Below is an overview of the key contributions made by each member:
\subsubsection{Project Manager and Backend Developer}
\textbf{Project Planning and Coordination:} Oversaw the project timeline, milestones, and deliverables, ensuring that all tasks were completed on schedule.
\textbf{Backend Development:} Responsible for designing the system’s database layout and coding inside the backend programming framework and API.
\textbf{Quality Assurance:} First screening was performed to check the fatal mistake and then fixed major issues that made the website bug-free and proper.


\subsubsection{Front-end Developer and QA Tester}
\textbf{User Interface Design:} Designed the layer located above where users interact with the system: it is most often clean and easily navigated and also responsive.

\noindent\textbf{Implementation of UI Components:} Developed and incorporated numerous user interface elements including elements that call for the recommended meals and elements that contain user profiles.

\noindent\textbf{Testing and Debugging:} Did functional testing, usability testing and also performance testing with a view of making sure the website delivers to the set specifications and guidelines.


\subsubsection{UI/UX Designer and Documentation Specialist}
\textbf{Design Concept and Wireframing:} Proposed the site architecture and sketched the first versions of the site, which helps in visualizing the website plan.

\noindent\textbf{User Experience Optimization:} Moreover, alternatively, carried out user and usability tests, which enabled to define the look and feel and optimize the user experience.

\noindent\textbf{Documentation:} Maintained comprehensive documentation of all major facets of the project, from designs and designs and programming structures to usage manuals.


\noindent\subsubsection{Collaborative Efforts}
In addition to individual contributions, the project team worked closely together in the following areas:
 
\noindent\textbf{Brainstorming and Ideation:} Specifically, the meetings presented in the text were designed to consist of team brainstorming during which ideas and effective solutions are shared.

\noindent\textbf{Code Reviews:} Continuously performed codewalk in order to keep up look-ups on the code and standard of the code to be adapted.

\noindent\textbf{Feedback and Iteration:} Offered feedback to one another throughout the project, finalizing the proposed solution and making the interface as effective and smooth as possible.


\subsubsection{External Contributions}
\textbf{Consultations with Nutrition Experts:} Consulted with nutritionists and dietitians to give assurance that all the nutritional information and advice as presented on the website were up-to-date and useful.

\vspace{12cm}

This site was a work in progress with contributions made by the whole team in consultation with our advisors. As demonstrated in the summary and outcomes of the project, the team has achieved the goal of creating an efficient, user-oriented platform to provide weight loss and diet recommendations tailored to individual needs.

% Bibliography
\clearpage
\bibliographystyle{ieeetr}
\addcontentsline{toc}{section}{References}
\printbibliography

\end{document}
